\section{Limitations and threats to validity}
\label{sec:limitations}

We state these rather than cash them in.

\paragraph{The \Sz{}/\So{} axis means different things for different method
classes.}
Two issues, in descending severity. First, \So{} is a \emph{no-op} for
model-free learned rows: they receive no parameters and run no forward model, so
their $\So/\Sz \approx 1.00$ carries no information about robustness and is never
tabled as such (Table~\ref{tab:classes}). Second, and applying even to the
\emph{conditioned} rows, being handed biased parameters is not the same exposure
as \emph{integrating} a biased model for a 500-step window: one is an input
perturbation, the other compounds through time. Experiment D2 exists to put both
on a common axis. Section~\ref{sec:marginal}'s intra-classical result needs
neither caveat and stands on its own.

\paragraph{``Your baselines were under-tuned.''}
The standard reply to any DA-versus-learned comparison. Our mitigation is
evidential rather than rhetorical: the localisation and inflation sensitivity
studies of Section~\ref{sec:tuning}, and the 4D-Var variance-scaling study that
returned a negative result. Both are reported as an appendix rather than left as
separate notes.

\paragraph{Different configurations must not be cross-quoted.}
Table~\ref{tab:price} and Table~\ref{tab:marginal} come from different
experimental configurations --- different observation counts, different
parameter-randomisation breadth. Ratios from one are not comparable to
absolute levels in the other, and we do not compare them.

\paragraph{The circularity objection.}
An amortised estimator needs a simulator to generate its training data. If that
simulator is the truth model, the comparison has handed it the truth --- and in
operations there is no truth model. This is the deepest objection to this line of
work and we do not claim to fully answer it. Our framing is therefore about
\emph{where prior information enters and what it costs}, not ``learned beats
physical''. The honest arm is a hybrid trained on a biased model and tested
against a different truth.
\todo{A self-supervised training objective --- a weak-constraint variational cost
evaluated on the estimator's own output, with no ground-truth supervision at all
--- is now implemented and is a direct, if partial, answer to this objection. It
has not been run. Decide whether it belongs in this paper or the companion.}

\paragraph{H2 as a strawman.}
No practitioner believes model fidelity is the only route to skill. We state H2
as \emph{binding}, not naive; D0 and D2 measure how weak the available buffers
are rather than assuming they are absent.

\paragraph{Idealised testbeds.}
Both systems are idealised. That is what makes ``matched everything'' definable
precisely enough for a marginal-value measurement to be meaningful, but a
referee will reasonably ask about transfer to operational systems. We own this
in scope rather than over-claim.

\section{Positioning}
\label{sec:positioning}

\todo{Every citation below is from memory and must be verified against the
literature before a bibliography exists. No reference in \texttt{refs.bib} has
been checked.}

The DA review literature frames the field's standard commitments; the
assimilation-in-the-unstable-subspace line gives the closest existing account of
\emph{why} DA's effective degrees of freedom are limited; weak-constraint 4D-Var
is the canonical relaxation of H2; innovation-based diagnostics supply the
standard consistency checks --- whose assumptions are violated when the estimator
is trained to fit the observations. On the learned side, the DA-with-learned-
model-error-correction line and the score-based / flow-based assimilation line
are the nearest neighbours.

\paragraph{The gap.}
The model-error literature asks how to \emph{correct} the model; the
unstable-subspace literature explains the dimensionality of the analysis; the
learned-DA literature proposes hybrids. As far as our survey goes, none of them
\emph{measures the marginal value of observations as a function of model error},
or attributes the deficit to named components of the inference operator.
Section~\ref{sec:marginal} is, on that reading, an unreported effect.
